\documentclass[11pt,a4paper]{article}
\usepackage[spanish,es-nodecimaldot]{babel}
\usepackage{fontspec}
\setmainfont{DejaVu Serif}\setsansfont{DejaVu Sans}
\usepackage[margin=1.9cm]{geometry}
\usepackage{array,tabularx,longtable,booktabs,enumitem,xcolor,hyperref,fancyhdr,lastpage}
\usepackage[most]{tcolorbox}
\definecolor{azul}{HTML}{224765}\definecolor{claro}{HTML}{F0F4F7}\definecolor{gris}{HTML}{4F5C68}
\hypersetup{colorlinks=true,linkcolor=azul,urlcolor=azul}
\setlist[itemize]{leftmargin=*,itemsep=2pt,topsep=3pt}
\setlist[enumerate]{leftmargin=*,itemsep=2pt,topsep=3pt}
\pagestyle{fancy}\fancyhf{}\fancyhead[L]{\small IA Total | Sesión 3}\fancyhead[R]{\small Guía de actividades}\fancyfoot[C]{\small \thepage\ de \pageref{LastPage}}\setlength{\headheight}{14pt}
\newtcolorbox{actividad}[1]{breakable,colback=claro,colframe=azul,title=\textbf{#1},fonttitle=\sffamily,boxrule=.5pt}
\newcommand{\diapos}[1]{\textcolor{gris}{\small\textsf{Diapositivas: #1}}\par}
\begin{document}
\begin{center}
{\Huge\sffamily\bfseries Plan de actividades: sesión 3}\\[6pt]
{\Large IA que crea: imágenes, audio y video}\\[5pt]
{\large Programa IA Total | San Juan Bosco, Morona Santiago}\\[8pt]
{\small Guía para coordinación entre instructor y asistente técnico | Versión 1.0 | Noviembre de 2026}
\end{center}
\begin{actividad}{Alcance del documento}
Esta guía indica qué actividades realizar, cuándo, con qué materiales y cómo participan los estudiantes. No establece las palabras del instructor ni sustituye sus explicaciones. La sesión trata sobre la IA que crea imágenes, audio y video: cómo funciona, cómo pedirle lo que queremos con un buen prompt, cómo reconocer contenido generado y cómo usarlo con responsabilidad para comunicar el proyecto final.
\end{actividad}

\section{Datos y objetivos}
\noindent\begin{tabularx}{\textwidth}{@{}p{3.3cm}X@{}}\toprule
\textbf{Grupo} & 40--50 estudiantes de bachillerato, cada uno con computador y celular conectados a internet; equipos de proyecto final de 4--5 integrantes.\\
\textbf{Duración} & 240 minutos presenciales, con una pausa de 15 minutos.\\
\textbf{Tema} & IA multimodal: generación de imágenes, texto a voz, voz a texto y video corto; reconocimiento de contenido generado; uso ético y propiedad intelectual.\\
\textbf{Materiales} & Presentación PDF (23 diapositivas); seis imágenes para «¿Real o generada?»; lista de cotejo (sección 4); acceso a Gemini, ChatGPT, Canva, Google Slides, TTSMaker y la aplicación CapCut del celular.\\
\textbf{Productos} & Por estudiante: imágenes de los niveles con sus prompts, una locución y una transcripción. Por equipo: portada, logo, locución y reel de 30 segundos del proyecto.\\\bottomrule
\end{tabularx}

\medskip
Al terminar, los estudiantes podrán explicar con sus palabras cómo una IA crea una imagen; escribir prompts de imagen con sujeto, escena, estilo, encuadre, luz y color, y lo que se quiere evitar; reconocer señales de contenido generado; convertir texto en voz y voz en texto revisando el resultado; armar un video corto con piezas generadas; y decidir cuándo es correcto publicar contenido hecho con IA.

\section{Cronograma general}
\small
\begin{longtable}{@{}p{2.0cm}p{1.1cm}p{5.2cm}p{5.6cm}@{}}
\toprule \textbf{Minutos} & \textbf{Dur.} & \textbf{Bloque} & \textbf{Modalidad}\\\midrule\endhead
00--05 & 5 & Inicio & Diapositivas 1--2. Objetivo y agenda.\\
05--30 & 25 & Avance del proyecto final & Diapositiva 3. Cada equipo, 2 minutos y una pregunta.\\
30--45 & 15 & ¿Real o generada? & Diapositivas 4--5. Todo el grupo, a mano alzada.\\
45--52 & 7 & Del ruido a la imagen & Diapositiva 6. Explicación y demostración.\\
52--65 & 13 & Fórmula del prompt de imagen & Diapositivas 7--8. Demostración construida con el grupo.\\
65--110 & 45 & Taller de imágenes & Diapositivas 9--14. Individual, en parejas y galería.\\
110--125 & 15 & Pausa & Instalar CapCut en el celular.\\
125--150 & 25 & Audio con IA & Diapositivas 15--17. Individual.\\
150--165 & 15 & ¿Está bien o no? & Diapositivas 18--19. Todo el grupo.\\
165--212 & 47 & Estudio del proyecto & Diapositivas 20--21. En equipos, con una pieza por integrante.\\
212--227 & 15 & Galería de proyectos & Diapositiva 22. Equipos.\\
227--240 & 13 & Cierre y tarea & Diapositiva 23.\\\bottomrule
\end{longtable}
\normalsize
La sesión abre con el avance de cada equipo y dedica la última hora a que cada equipo cree las piezas para comunicar su proyecto: portada, logo, locución y reel.

\section{Desarrollo de la sesión}
\subsection{00--05 | Inicio}
\diapos{1--2.}
Presentar el objetivo (crear imágenes, audio y video con IA, y saber cuándo publicarlos) y la agenda. Los equipos se sientan juntos.

\subsection{05--30 | Avance del proyecto final}
\diapos{3.}
Cada equipo cuenta en 2 minutos qué averiguó desde la sesión anterior y cómo va su idea. Al terminar, el equipo siguiente le hace una pregunta. Con 8--12 equipos son 20--25 minutos. El asistente controla el tiempo; el instructor anota el estado de cada proyecto (idea clara, en definición o sin avance) y al final da una retroalimentación general breve.

\subsection{30--45 | ¿Real o generada?}
\diapos{4--5.}
\begin{enumerate}
\item 30--40 (diapositiva 4). Se proyectan seis imágenes, tres fotografías reales de San Juan Bosco o la Amazonía y tres generadas con IA, sin decir cuáles son. Para cada número, mano arriba quien crea que la hizo una IA. Antes de revelar la respuesta se pregunta a dos estudiantes en qué se fijaron.
\item 40--45 (diapositiva 5). Repasar las pistas: manos y dedos, texto deforme, sombras o reflejos sin sentido y aspecto demasiado perfecto. Cerrar con la idea de que cada vez es más difícil notarlo, por eso importa saber de dónde viene una imagen.
\end{enumerate}

\subsection{45--52 | Del ruido a la imagen}
\diapos{6.}
Un generador de imágenes empieza con puntos al azar, como la estática de una televisión, y paso a paso los ordena hasta que la imagen se parece a lo que se pidió. Igual que con el texto, la IA no busca una imagen «verdadera»: produce algo probable según lo que aprendió de millones de imágenes con sus descripciones. El instructor genera dos veces el mismo prompt y el grupo compara las dos imágenes.

\subsection{52--65 | Fórmula del prompt de imagen}
\diapos{7--8.}
\begin{enumerate}
\item 52--56 (diapositiva 7). Presentar la fórmula: sujeto, escena, estilo, encuadre, luz y color, y lo que se quiere evitar. Los estudiantes señalan cada elemento en el ejemplo del tucán. Recordar la fórmula de la sesión anterior: rol y tono no hacen falta en una imagen, pero sí ser específico y poner límites.
\item 56--65 (diapositiva 8). El instructor genera «Un paisaje de la Amazonía» y, en tres rondas, el grupo dicta qué agregar: qué aparece y dónde, en qué estilo, y qué encuadre, luz y colores. Se genera después de cada ronda y se comparan los resultados.
\end{enumerate}

\subsection{65--110 | Taller de imágenes}
\diapos{9--14.}
Cada estudiante trabaja en su computador con Gemini y ChatGPT; si una llega a su límite diario, sigue con la otra. Para el nivel 3 entran a Canva, con su cuenta de Google si no tienen otra. Guardan cada imagen con su prompt, porque los usarán en el proyecto. Al terminar cada nivel, el instructor pide a uno o dos estudiantes que muestren su imagen y lean su prompt. El asistente marca la lista de cotejo mientras recorre el aula.
\begin{enumerate}
\item 65--67 (diapositiva 9): explicar cómo funciona el taller.
\item 67--75, nivel 1 (diapositiva 10): el mismo prompt en Gemini y en ChatGPT; comparar cuál se parece más a lo pedido y si respetó lo que se quería evitar.
\item 75--83, nivel 2 (diapositiva 11): un sujeto de San Juan Bosco en tres estilos (fotografía, acuarela e ilustración plana), cambiando solo el estilo.
\item 83--95, nivel 3 (diapositiva 12): una imagen vertical para redes que invite a jóvenes de otra ciudad a visitar San Juan Bosco. Se genera «sin texto» y el texto se agrega en Canva con una plantilla de historia de Instagram, porque la IA suele escribir mal. Se revisa si muestra lugares reales o inventados y se añade la etiqueta «Hecho con IA».
\item 95--105, nivel 4 (diapositiva 13), teléfono de imágenes en parejas: uno genera una imagen sin mostrarla y la describe en voz alta con la fórmula; el otro escribe el prompt y genera la suya. Comparan qué detalle faltó y cambian de papel.
\item 105--110, galería caminante (diapositiva 14): cada estudiante deja en pantalla su mejor imagen y el prompt que usó; todos recorren el aula. Dos o tres cuentan qué prompt les gustaría copiar y qué imagen se ve más real.
\end{enumerate}

\subsection{110--125 | Pausa}
Antes de salir, cada estudiante comprueba que tiene la aplicación CapCut en el celular. Se usa sin iniciar sesión: basta con abrirla y aceptar los permisos.

\subsection{125--150 | Audio con IA}
\diapos{15--17.}
\begin{enumerate}
\item 125--127 (diapositiva 15): tres usos del audio con IA: texto a voz, voz a texto y clonación de voz. Hoy se practican los dos primeros; la clonación solo se menciona y se retoma en «¿Está bien o no?».
\item 127--139, texto a voz (diapositiva 16). Cada estudiante pide a un asistente el guion de una invitación de 50 palabras a la Feria Comunitaria con el prompt de la diapositiva, lo corrige, lo pega en TTSMaker (ttsmaker.com, funciona sin cuenta), prueba dos voces en español y descarga el MP3. Comentan si suena natural y qué palabras pronuncia mal. Guardan el audio.
\item 139--150, voz a texto (diapositiva 17). Graban una nota de voz de 20 segundos en el celular, suben el audio a Gemini y le piden la transcripción; revisan qué palabras escribió mal. Quien hable kichwa o shuar prueba también en su idioma y el grupo compara los resultados.
\end{enumerate}

\subsection{150--165 | ¿Está bien o no?}
\diapos{18--19.}
\begin{enumerate}
\item 150--162 (diapositiva 18). El instructor lee cada caso y los estudiantes responden con la mano: pulgar arriba (está bien), mano plana (depende) o pulgar abajo (no está bien). Después de cada caso, uno o dos explican su respuesta.
  \begin{itemize}
  \item Imagen generada para el afiche de la Feria, diciendo que es IA: está bien.
  \item Clonar la voz de un profesor para una broma: no está bien; es usar la voz de alguien sin permiso.
  \item Foto «real» de un compañero en un lugar donde nunca estuvo: no está bien; así se hacen los deepfakes.
  \item Imagen «al estilo de» un artista famoso para vender camisetas: depende, y casi siempre no; el estilo y el trabajo del artista tienen dueño.
  \item Video hecho con IA en redes sin decir que lo es: depende del contenido, pero lo correcto es indicarlo.
  \end{itemize}
\item 162--165 (diapositiva 19): cerrar con las cuatro reglas antes de publicar: decir que se hizo con IA, no usar la cara ni la voz de alguien sin permiso, no copiar marcas ni el estilo de un artista para vender, y revisar que no muestre lugares o datos falsos.
\end{enumerate}

\subsection{165--212 | Estudio del proyecto}
\diapos{20--21.}
Cada equipo crea las piezas para comunicar su proyecto. Cada integrante trabaja en su computador con una pieza:
\begin{itemize}
\item Portada: imagen principal del proyecto, con la fórmula de imagen.
\item Logo: un ícono simple y el nombre del proyecto en Google Slides.
\item Locución: guion de 20 segundos que explique el proyecto, convertido en voz con TTSMaker. El MP3 se pasa al celular de quien arma el reel.
\item Reel: junta las piezas en la aplicación CapCut del celular, sin iniciar sesión.
\item Revisor: comprueba las reglas de «Antes de publicar» en cada pieza. En equipos de cuatro, también arma el reel.
\end{itemize}
\begin{enumerate}
\item 165--168: reparto de piezas y acuerdo sobre estilo y colores comunes.
\item 168--190: cada integrante crea su pieza y la comparte con quien arma el reel.
\item 190--212 (diapositiva 21): el reel de 30 segundos con portada, logo y dos imágenes más, la locución, subtítulos automáticos corregidos y un cierre con el nombre del proyecto y «Hecho con IA». Se exporta y se guarda para subirlo a Moodle.
\end{enumerate}
El instructor y el asistente recorren los equipos y ayudan con CapCut.

\subsection{212--227 | Galería de proyectos}
\diapos{22.}
Cada equipo deja su reel listo en una pantalla y los estudiantes recorren los puestos. Cada equipo recibe de los demás una estrella (algo que funciona) y un deseo (algo que mejorar). El instructor comenta dos o tres reels al final.

\subsection{227--240 | Cierre y tarea}
\diapos{23.}
\begin{enumerate}
\item 227--233, ronda final: cada equipo dice qué pieza le gustó más y qué cuidaría al publicarla.
\item 233--240, tarea en Moodle: individual, subir su mejor imagen con el prompt y responder dos preguntas sobre uso ético; en equipo, subir en «Avance 3» la portada, el logo, la locución y el reel del proyecto.
\end{enumerate}

\subsection{Ajuste del tiempo}
Si el grupo avanza rápido, se alargan los niveles 3 y 4 del taller y el estudio del proyecto, y se puede repetir la galería caminante con las imágenes del nivel 4. Si avanza lento, se omite el nivel 4, la voz a texto se hace como demostración del instructor y la galería de proyectos se reduce a cuatro equipos que muestran su reel en el proyector.

\section{Evaluación}
La evaluación usa tres fuentes:
\begin{enumerate}
\item La lista de cotejo, que el asistente completa mientras recorre el aula. Cada casilla se marca con {\sffamily ✓} (logrado), \textasciitilde\ (en proceso) o se deja vacía.
\item Lo que cada estudiante produce: imágenes con sus prompts, locución y transcripción; en Moodle, su mejor imagen y la reflexión sobre uso ético.
\item Lo que produce cada equipo: portada, logo, locución y reel del proyecto.
\end{enumerate}

\medskip
\footnotesize
\noindent\begin{tabularx}{\textwidth}{@{}l*{7}{>{\centering\arraybackslash}X}@{}}\toprule
\textbf{Estudiante} & \textbf{Nivel 1--2} & \textbf{Nivel 3} & \textbf{Teléfono} & \textbf{Locución} & \textbf{Trans\-cripción} & \textbf{Pieza} & \textbf{Reglas}\\\midrule
\dots & & & & & & & \\\bottomrule
\end{tabularx}

\smallskip
\textcolor{gris}{Nivel 1--2: comparó herramientas y estilos. Nivel 3: imagen vertical con texto agregado en Canva y etiqueta «Hecho con IA». Teléfono: describió y recreó una imagen en pareja. Locución: generó y escuchó la invitación. Transcripción: revisó los errores. Pieza: terminó su pieza del proyecto. Reglas: su pieza cumple las reglas de «Antes de publicar».}

\medskip
\noindent\begin{tabularx}{\textwidth}{@{}p{2.6cm}XXX@{}}\toprule
\textbf{Criterio} & \textbf{Inicial} & \textbf{En proceso} & \textbf{Logrado}\\\midrule
Descripción & El prompt solo nombra el sujeto. & Incluye sujeto y escena. & Incluye sujeto, escena, estilo, encuadre y luz o color.\\
Control & No revisa la imagen. & Revisa si se parece a lo pedido. & Ajusta el prompt hasta lograr lo que quería y usa «evita» cuando hace falta.\\
Uso responsable & No indica que usó IA. & Indica que usó IA. & Indica que usó IA y no usa personas, marcas ni lugares falsos.\\\bottomrule
\end{tabularx}
\normalsize

\section{Recursos y preparación previa}
\begin{itemize}
\item Presentación: PDF de 23 diapositivas, proyector y cable o adaptador probados.
\item «¿Real o generada?»: seis imágenes en \texttt{sesion3/presentation/images/}, recortadas en proporción 8:5. Tres fotografías reales de San Juan Bosco o la Amazonía, tomadas por el equipo o con permiso de uso, y tres generadas con los prompts del archivo de instrucciones de imágenes de la sesión. Anotar cuáles son reales.
\item Herramientas: comprobar en la red del aula Gemini y ChatGPT (generación de imágenes y transcripción de audio), Canva, Google Slides, TTSMaker (texto a voz sin cuenta, descarga en MP3), la aplicación CapCut del celular sin iniciar sesión (edición, subtítulos automáticos y exportación).
\item Respuestas de «¿Real o generada?»: las imágenes 1, 3 y 6 son fotografías reales; las 2, 4 y 5 son generadas (2: demasiado perfecta; 4: manos; 5: texto del letrero).
\item Asistente técnico: lista de cotejo con los nombres de los estudiantes agrupados por equipo.
\end{itemize}

\section{Responsabilidades}
\noindent\begin{tabularx}{\dimexpr\textwidth-4pt\relax}{@{}p{3cm}X@{}}\toprule
\textbf{Instructor} & Dirige «¿Real o generada?», las demostraciones y los casos de «¿Está bien o no?»; comenta los reels.\\
\textbf{Asistente técnico} & Resuelve problemas de acceso, controla los tiempos, recorre el aula, completa la lista de cotejo, apoya con CapCut y anota el estado de los proyectos.\\
\textbf{Estudiantes} & Crean y comparan imágenes, describen y recrean imágenes en pareja, generan y transcriben audio, opinan sobre los casos y crean una pieza del proyecto con su equipo.\\\bottomrule
\end{tabularx}

\section{Contingencias}
\begin{itemize}
\item Una herramienta de imágenes llega a su límite diario: cambiar entre Gemini y ChatGPT.
\item TTSMaker no carga: escuchar el guion con el botón de audio de Google Traductor o con «Leer en voz alta» de Microsoft Edge y grabarlo con el celular.
\item CapCut pide iniciar sesión para alguna función: usar solo la edición básica (imágenes, audio, texto y subtítulos); si no deja exportar, armar el reel en el celular de otro integrante.
\item Quien no tenga la edad que exige una herramienta trabaja en el computador de un compañero.
\item Ritmo más rápido o más lento de lo previsto: ver «Ajuste del tiempo».
\item Grupo distraído: acortar la actividad en curso, pasar al momento de compartir y pedir a quien se distrae que muestre su pantalla.
\item Privacidad y derechos: no subir fotos de compañeros ni de otras personas a las herramientas, no clonar voces de otras personas, no usar logotipos de marcas y etiquetar con «Hecho con IA» todo lo que se publique.
\end{itemize}
\end{document}
